\usepackage{setspace}
\usepackage{graphicx} %绘图
\usepackage{pifont} %圆圈1，如①
\usepackage{booktabs} %绘制三线表
\usepackage{multirow} %使用https://www.tablesgenerator.com/网站绘制表格所需包
\usepackage{threeparttable} %表格脚注
\usepackage[noend]{algpseudocode} %绘制伪代码
\usepackage{algorithmicx,algorithm}
\usepackage{pythonhighlight} %python代码高亮

% ===== 代码样式配置 =====
\lstdefinestyle{cstyle}{
  language=C,
  basicstyle=\small\ttfamily,
  keywordstyle=\color{blue}\bfseries,
  commentstyle=\color{green!50!black}\itshape,
  stringstyle=\color{red!60!black},
  numberstyle=\tiny\color{gray},
  numbers=left,
  stepnumber=1,
  numbersep=8pt,
  showstringspaces=false,
  breaklines=true,
  breakatwhitespace=false,
  frame=single,
  framerule=0.5pt,
  rulecolor=\color{gray!40},
  backgroundcolor=\color{gray!5},
  tabsize=4,
  captionpos=b,
  xleftmargin=2em,
  xrightmargin=0.5em,
  aboveskip=0.8em,
  belowskip=0.8em,
}
\lstdefinestyle{pythonstyle}{
  language=Python,
  basicstyle=\small\ttfamily,
  keywordstyle=\color{blue}\bfseries,
  commentstyle=\color{green!50!black}\itshape,
  stringstyle=\color{red!60!black},
  numberstyle=\tiny\color{gray},
  numbers=left,
  stepnumber=1,
  numbersep=8pt,
  showstringspaces=false,
  breaklines=true,
  frame=single,
  framerule=0.5pt,
  rulecolor=\color{gray!40},
  backgroundcolor=\color{gray!5},
  tabsize=4,
  captionpos=b,
  xleftmargin=2em,
  xrightmargin=0.5em,
  aboveskip=0.8em,
  belowskip=0.8em,
}
\renewcommand{\lstlistingname}{代码}
\usepackage[unicode=true,colorlinks=true,linkcolor=black,anchorcolor=black,citecolor=black]{hyperref}%引用参考文献、章节、图表，并跳转
\hypersetup{bookmarksnumbered=true,bookmarksopen=true,citecolor=black}
\AtBeginDocument{\hypersetup{citecolor=black}}
\usepackage{caption}
\captionsetup[figure]{font=small,labelsep=space} 
\captionsetup[table]{font=small,labelsep=space} %图片和表格修改caption

\usepackage[final]{pdfpages}
% 设置页边距
\usepackage{geometry}
\usepackage{makecell}
\usepackage{caption}
\usepackage{multirow}
\usepackage{booktabs}
\makeatletter
\let\c@lofdepth\relax
\let\c@lotdepth\relax
\makeatother
\usepackage{subfigure}
\usepackage{float}
\usepackage{placeins}
\usepackage[titles,subfigure]{tocloft}
\usepackage{soul}
\usepackage{color,xcolor}



\geometry{a4paper,top=3cm,bottom=2.5cm,left=2.5cm,right=2.5cm}

% % 设置行间距 1.5倍
% \onehalfspacing

\linespread{1.5}\selectfont
% 设置段与段之间的垂直距离 \parskip默认橡皮长度是0pt plus 1pt
\setlength{\parskip}{0pt}
% 正文首行缩进，符合学校模板要求
\setlength{\parindent}{2em}
\emergencystretch=1.5em
% 允许页面底部不强制对齐，减少 underfull vbox 警告
\raggedbottom
% \setlength{\parindent}{0pt}

%设置字体
% 设置英文字体
\setmainfont{Times New Roman}[
    BoldFont = Times New Roman Bold,
    ItalicFont = Times New Roman Italic,
    BoldItalicFont = Times New Roman Bold Italic
]

\setsansfont{Times New Roman}[
    BoldFont = Times New Roman Bold,
    ItalicFont = Times New Roman Italic,
    BoldItalicFont = Times New Roman Bold Italic
]

\setmonofont{Times New Roman}[
    BoldFont = Times New Roman Bold,
    ItalicFont = Times New Roman Italic,
    BoldItalicFont = Times New Roman Bold Italic
]

%设置字号
\usepackage{ctexsize,type1cm}
\newcommand{\yihao}{\fontsize{26pt}{39pt}\selectfont}
\newcommand{\xiaoyi}{\fontsize{24pt}{36pt}\selectfont}   
\newcommand{\erhao}{\fontsize{22pt}{33pt}\selectfont}          
\newcommand{\xiaoer}{\fontsize{18pt}{27pt}\selectfont}          
\newcommand{\sanhao}{\fontsize{16pt}{24pt}\selectfont}        
\newcommand{\xiaosan}{\fontsize{15pt}{22.5pt}\selectfont}        
\newcommand{\sihao}{\fontsize{14pt}{21pt}\selectfont}            
\newcommand{\xiaosi}{\fontsize{12pt}{18pt}\selectfont}            
\newcommand{\wuhao}{\fontsize{10.5pt}{15.75pt}\selectfont}
\newcommand{\xiaowu}{\fontsize{9pt}{13.5pt}\selectfont}    
\newcommand{\liuhao}{\fontsize{7.5pt}{11.25pt}\selectfont}

%使用公式，表格，图片
\usepackage{mathtools,amsmath,amssymb,graphicx,array,float}

% 设置页眉面脚
%% 设置章节前的页码格式
\usepackage{fancyhdr}
\fancypagestyle{content}{
	\fancyhf{}
	\renewcommand{\headrulewidth}{0pt}
	\fancyfoot[C]{\songti\xiaowu  \Roman{page}}
}

%重新设置headings
\fancypagestyle{headings}{
    \fancyhf{}
	\fancyfoot[C]{\songti\xiaowu  \Roman{page} }
	\renewcommand{\headrulewidth}{0.5pt}
	\fancyhead[CO]{\songti\xiaowu 目录 }
	\fancyhead[CE]{\songti\xiaowu 重庆邮电大学本科毕业设计（论文）}
}

%重新设置plain
%正文页眉设置
\fancypagestyle{plain}{
    \fancyhf{}
    \fancyfoot[C]{\songti\xiaowu  \thepage }
    \fancyhead[CO]{\songti\xiaowu \leftmark }
    \fancyhead[CE]{\songti\xiaowu 重庆邮电大学本科毕业设计（论文）}
}


%设置双线页脚
\makeatletter
\def\footrule{
    {\if@fancyplain\let\footrulewidth\plainfootrulewidth\fi%
    \hrule\@height 0pt \@width\headwidth          %上面0.5pt粗
    \vskip 1pt
    \hrule\@height 0pt \@width\headwidth %下面线为1pt粗
    \vskip-2\headrulewidth\vskip-1.2pt}    %两条线的距离1pt
    \vspace{2mm}}     %双线与下面正文之间的垂直间距
\makeatother

%设置文章格式
\ctexset {
    contentsname={目录},
    listfigurename={插图},
    listtablename={表格},
    figurename={图},
    tablename={表},
    bibname={参考文献},
    appendixname={附录},
    chapter={
        beforeskip={-17.5pt},
        afterskip={16.5pt},
        nameformat={\heiti\sanhao\centering},
        number={\arabic{chapter}},
        titleformat={\heiti\sanhao\centering},
    },
    section={
        format={\heiti\sihao},
        beforeskip={13pt},
        afterskip={13pt},
    },
    subsection={
        format={\heiti\xiaosi},
        beforeskip={13pt},
        afterskip={13pt},
    },
}

% 目录中的章加点
\usepackage[titles]{tocloft}
\renewcommand{\cftdot}{$\cdot$}
\renewcommand{\cftdotsep}{1.5}
\setlength{\cftbeforechapskip}{0.5em}
\setlength{\cftbeforesecskip}{0.5em}
\setlength{\cftbeforesubsecskip}{0.5em}

\renewcommand{\cftchapleader}{\cftdotfill{\cftchapdotsep}}
\renewcommand{\cftchapdotsep}{\cftdotsep}
\makeatletter
\renewcommand{\numberline}[1]{%
\settowidth\@tempdimb{#1\hspace{0.5em}}%
\ifdim\@tempdima<\@tempdimb%
  \@tempdima=\@tempdimb%
\fi%
\hb@xt@\@tempdima{\@cftbsnum #1\@cftasnum\hfil}\@cftasnumb}
\makeatother

% 设置目录字体尺寸
\renewcommand{\cftchapfont}{\songti\xiaosi\mdseries}
\renewcommand{\cftchappagefont}{\songti\xiaosi\mdseries}
\renewcommand{\cftsecfont}{\songti\xiaosi\mdseries}
\renewcommand{\cftsecpagefont}{\songti\xiaosi\mdseries}
\renewcommand{\cftsubsecfont}{\songti\xiaosi\mdseries}
\renewcommand{\cftsubsecpagefont}{\songti\xiaosi\mdseries}

%使用代码排版包
\usepackage{listings}
\usepackage{color}
\lstset{%
    frame=shadowbox,
    extendedchars=false,            % 不使用xelatex而使用CJK方式处理汉字
    language=python,
    basicstyle=\sffamily,           % 设置整体格式
    keywordstyle=\bfseries,         % 关键字格式
    commentstyle=\rmfamily\itshape, % 注释格式
    stringstyle=\ttfamily,          % 字符串格式
    columns=flexible,
    escapechar=',                   % 注释中显示汉字，eg //'一个整数'
    tabsize=4,
    numbers=left,
    numberstyle=\small,             % 行号字体设置
    stepnumber=1,                   % 行号距离设置，1代表每行加行号
    numbersep=8pt,                  % 行号和代码距离设置
    backgroundcolor=\color{white},
    showspaces=false,               % show spaces adding particular underscores
    showstringspaces=false,         % 使用下划线连接字符串
    showtabs=false,
    frame=single,                   % 给代码加边框
    captionpos=b,                   % sets the caption-position to bottom
    breaklines=true,                % 自动换行设置
    breakatwhitespace=false,        % sets if automatic breaks should only happen at whitespace
    escapeinside={\%*}{*)},         % if you want to add a comment within your code
    xleftmargin=2em,                % 设置左边距，宽度默认是与页芯等宽的
    xrightmargin=2em,               % 设置右边距，宽度默认是与页芯等宽的
    aboveskip=1em                   % 设置上边距
}

%设置自定义变量
\newcommand\degree{^\cire}

% 定义文献引用格式：正文中默认使用右上角角标引用
\usepackage{cite}
\let\citeinline\cite
\renewcommand{\cite}[1]{\textsuperscript{\citeinline{#1}}}
\newcommand{\upcite}[1]{\cite{#1}}
\newcommand\supercite[2][]{%
\textsuperscript{\citeinline[#1]{#2}}
}

\usepackage{enumitem}
\setlist[description]{
    itemsep=-5pt,
    font=\songti,
}

% 定义中文封面环境
\newenvironment{titletabbing}
{\par\bfseries\songti\sihao\tabbing}
{\endtabbing\par}

\usepackage[nottoc]{tocbibind}
\endinput
